\QHeading{第10题　口算。}
\par\noindent


\begin{oralgrid}
\OralItem{1}{240\div30}{8}
\OralItem{2}{19\times4}{76}
\OralItem{3}{35\times20}{700}
\OralItem{4}{630\div90}{7}
\OralItem{5}{12\times8}{96}
\OralItem{6}{50\times16}{800}
\OralItem{7}{720\div9}{80}
\OralItem{8}{26\times3}{78}
\OralItem{9}{45+55}{100}
\OralItem{10}{130-70}{60}
\OralItem{11}{320\div8}{40}
\OralItem{12}{15\times3}{45}\end{oralgrid}
\par

\QHeading{口算思路补充（对应上表题号）}

\Answer{上面列出得数，下面按小题给出不同计算思路。计数单位法适用于小数，整数题可比较凑整、拆分或商不变等方法。}
\MethodTitle{第1小题：$240\div 30$}

\MethodBlock{口算思路一：同时去掉相同个数的0}{}{\begin{aligned}[t]&\phantom{={}}240\div 30\\
&=(240\div 10)\div (30\div 10)\\
&=24\div 3\\
&=8\end{aligned}}

\MethodBlock{口算思路二：拆分除数，连续相除}{}{\begin{aligned}[t]&\phantom{={}}240\div 30\\
&=240\div (2\times 15)\\
&=240\div 2\div 15\\
&=120\div 15\\
&=8\end{aligned}}
\MethodTitle{第2小题：$19\times 4$}

\MethodBlock{口算思路一：凑整后补偿}{把接近整十、整百的因数凑整，再补回或减去多算的部分。}{\begin{aligned}[t]&\phantom{={}}19\times 4\\
&=4\times (20-1)\\
&=4\times 20-4\times 1\\
&=80-4\\
&=76\end{aligned}}

\MethodBlock{口算思路二：按数位拆分计算}{把一个因数拆成整十、整百数与余下部分之和。}{\begin{aligned}[t]&\phantom{={}}19\times 4\\
&=4\times (10+9)\\
&=4\times 10+4\times 9\\
&=40+36\\
&=76\end{aligned}}
\MethodTitle{第3小题：$35\times 20$}

\MethodBlock{口算思路一：按数位拆分计算}{把一个因数拆成整十、整百数与余下部分之和。}{\begin{aligned}[t]&\phantom{={}}35\times 20\\
&=20\times (30+5)\\
&=20\times 30+20\times 5\\
&=600+100\\
&=700\end{aligned}}
\MethodTitle{第4小题：$630\div 90$}

\MethodBlock{口算思路一：同时去掉相同个数的0}{}{\begin{aligned}[t]&\phantom{={}}630\div 90\\
&=(630\div 10)\div (90\div 10)\\
&=63\div 9\\
&=7\end{aligned}}

\MethodBlock{口算思路二：拆分除数，连续相除}{}{\begin{aligned}[t]&\phantom{={}}630\div 90\\
&=630\div (2\times 45)\\
&=630\div 2\div 45\\
&=315\div 45\\
&=7\end{aligned}}
\MethodTitle{第5小题：$12\times 8$}

\MethodBlock{口算思路一：凑整后补偿}{把接近整十、整百的因数凑整，再补回或减去多算的部分。}{\begin{aligned}[t]&\phantom{={}}12\times 8\\
&=12\times (10-2)\\
&=12\times 10-12\times 2\\
&=120-24\\
&=96\end{aligned}}
\MethodTitle{第6小题：$50\times 16$}

\MethodBlock{口算思路一：拆因数，凑整计算}{把一个因数拆成两个因数，先凑成整百或整千。}{\begin{aligned}[t]&\phantom{={}}50\times 16\\
&=50\times (2\times 8)\\
&=(50\times 2)\times 8\\
&=100\times 8\\
&=800\end{aligned}}

\MethodBlock{口算思路二：按数位拆分计算}{把一个因数拆成整十、整百数与余下部分之和。}{\begin{aligned}[t]&\phantom{={}}50\times 16\\
&=50\times (10+6)\\
&=50\times 10+50\times 6\\
&=500+300\\
&=800\end{aligned}}
\MethodTitle{第7小题：$720\div 9$}

\MethodBlock{口算思路一：拆分除数，连续相除}{}{\begin{aligned}[t]&\phantom{={}}720\div 9\\
&=720\div (3\times 3)\\
&=720\div 3\div 3\\
&=240\div 3\\
&=80\end{aligned}}
\MethodTitle{第8小题：$26\times 3$}

\MethodBlock{口算思路一：按数位拆分计算}{把一个因数拆成整十、整百数与余下部分之和。}{\begin{aligned}[t]&\phantom{={}}26\times 3\\
&=3\times (20+6)\\
&=3\times 20+3\times 6\\
&=60+18\\
&=78\end{aligned}}
\MethodTitle{第10小题：$130-70$}

\MethodBlock{口算思路一：补差法}{从减数先补到一个整数，再补到被减数，两段差相加。}{\begin{aligned}[t]&\phantom{={}}130-70\\
&=(80-70)+(130-80)\\
&=10+50\\
&=60\end{aligned}}
\MethodTitle{第11小题：$320\div 8$}

\MethodBlock{口算思路一：利用商不变的性质}{被除数和除数同时乘同一个非零数，商不变。}{\begin{aligned}[t]&\phantom{={}}320\div 8\\
&=(320\times 5)\div (8\times 5)\\
&=1600\div 40\\
&=40\end{aligned}}

\MethodBlock{口算思路二：拆分除数，连续相除}{}{\begin{aligned}[t]&\phantom{={}}320\div 8\\
&=320\div (2\times 4)\\
&=320\div 2\div 4\\
&=160\div 4\\
&=40\end{aligned}}
\MethodTitle{第12小题：$15\times 3$}

\MethodBlock{口算思路一：按数位拆分计算}{把一个因数拆成整十、整百数与余下部分之和。}{\begin{aligned}[t]&\phantom{={}}15\times 3\\
&=3\times (10+5)\\
&=3\times 10+3\times 5\\
&=30+15\\
&=45\end{aligned}}


\par\noindent

\QHeading{第11题　脱式计算。}

\CalcBlock{1}{\begin{aligned}[t]&\phantom{={}}\color{black}32+(76-24)\div4\\&\color{answerred}=32+52\div4\\&\color{answerred}=32+13\\&\color{answerred}=45\end{aligned}}
\par

\CalcBlock[方法一]{2}{\begin{aligned}[t]&\phantom{={}}\color{black}(45+15)\times(28\div7)\\&\color{answerred}=60\times4\\&\color{answerred}=240\end{aligned}}
\MethodBlock{方法二：乘法分配律}{分别求出两部分的积，再合起来。}{\begin{aligned}[t]&\phantom{={}}(45+15)\times (28\div 7)\\
&=45\times (28\div 7)+15\times (28\div 7)\\
&=45\times 4+15\times 4\\
&=180+60\\
&=240\end{aligned}}

\CalcBlock[方法一]{3}{\begin{aligned}[t]&\phantom{={}}\color{black}180-(36+48\div6)\\&\color{answerred}=180-(36+8)\\&\color{answerred}=180-44\\&\color{answerred}=136\end{aligned}}
\MethodBlock{方法二：先处理括号或乘除，再去括号后计算}{加号后的括号去掉不变号，减号后的括号去掉要改变各项的加减号。}{\begin{aligned}[t]&\phantom{={}}180-(36+48\div 6)\\
&=180-(36+8)\\
&=180-36-8\\
&=144-8\\
&=136\end{aligned}}

\CalcBlock{4}{\begin{aligned}[t]&\phantom{={}}\color{black}72\div(12-4\times2)\\&\color{answerred}=72\div(12-8)\\&\color{answerred}=72\div4\\&\color{answerred}=18\end{aligned}}
\CalcBlock{5}{\begin{aligned}[t]&\phantom{={}}\color{black}25\times(30-12)-180\\&\color{answerred}=25\times18-180\\&\color{answerred}=450-180\\&\color{answerred}=270\end{aligned}}
\CalcBlock[方法一]{6}{\begin{aligned}[t]&\phantom{={}}\color{black}[64+(90-58)]\times2\\&\color{answerred}=[64+32]\times2\\&\color{answerred}=96\times2\\&\color{answerred}=192\end{aligned}}
\MethodBlock{方法二：乘法分配律}{分别求出两部分的积，再合起来。}{\begin{aligned}[t]&\phantom{={}}[64+(90-58)]\times 2\\
&=64\times 2+(90-58)\times 2\\
&=64\times 2+32\times 2\\
&=128+64\\
&=192\end{aligned}}

\par\noindent

\QHeading{第12题　解决问题。}

\Problem{①用一辆载质量为6吨的卡车运送大米和面粉，大米每袋50千克，面粉每袋45千克。在装了60袋面粉后，这辆卡车还能装多少袋大米？}
\par


\MethodTitle{方法一}
\Answer{$6\text{吨}=6000\text{千克}$。\par \Formula{\begin{aligned}[t]&\phantom{={}}(6000-45\times60)\div50\\&=(6000-2700)\div50\\&=3300\div50\\&=66\text{（袋）}\end{aligned}}\par 答：还能装66袋大米。}

\MethodTitle{方法二：把载质量都折合成大米的袋数}
\Answer{6吨是6000千克。先求只装大米能装多少袋，再扣去面粉占用的相同载质量。}
\Answer{\Formula{\begin{aligned}[t]&\phantom{={}}6000\div 50-(45\times 60)\div 50\\
&=6000\div 50-2700\div 50\\
&=120-54\\
&=66\text{（袋）}\end{aligned}}}
\Answer{答：还能装66袋大米。}
\Problem{②李老师去文具店买文具，买了8支钢笔和12本笔记本，一共用去156元。每本笔记本5元，每支钢笔多少钱？}

\MethodTitle{方法一}
\Answer{\Formula{\begin{aligned}[t]&\phantom{={}}(156-5\times12)\div8\\&=(156-60)\div8\\&=96\div8\\&=12\text{（元）}\end{aligned}}\par 答：每支钢笔12元。}

\MethodTitle{方法二：平均分成4组购买}
\Answer{8支钢笔和12本笔记本平均分成4组，每组2支钢笔、3本笔记本，共39元。}
\Answer{\textbf{每组总价}\par \Formula{\begin{aligned}[t]&\phantom{={}}156\div 4\\
&=39\text{（元）}\end{aligned}}}
\Answer{\textbf{每组钢笔的价钱}\par \Formula{\begin{aligned}[t]&\phantom{={}}39-5\times 3\\
&=39-15\\
&=24\text{（元）}\end{aligned}}}
\Answer{\textbf{每支钢笔的价钱}\par \Formula{\begin{aligned}[t]&\phantom{={}}24\div 2\\
&=12\text{（元）}\end{aligned}}}
\Answer{答：每支钢笔12元。}
